% Generated draft: requires mathematical and source review.
\documentclass[11pt]{article}
\usepackage[margin=25mm]{geometry}
\usepackage{fontspec}
\setmainfont{Latin Modern Roman}
\usepackage{amsmath,amssymb,mathtools,unicode-math}
\setmathfont{Latin Modern Math}
\usepackage{enumitem,needspace,xurl,hyperref}
\setlength{\parindent}{0pt}
\setlength{\parskip}{5pt}
\setlength{\emergencystretch}{3em}
\newcommand{\N}{\mathbb{N}}
\newcommand{\Z}{\mathbb{Z}}
\newcommand{\Q}{\mathbb{Q}}
\newcommand{\R}{\mathbb{R}}
\newcommand{\C}{\mathbb{C}}
\newcommand{\F}{\mathbb{F}}
\newcommand{\A}{\mathbb{A}}
\newcommand{\PP}{\mathbb{P}}
\newcommand{\E}{\mathbb{E}}
\newcommand{\eps}{\varepsilon}
\newcommand{\dd}{\,\mathrm{d}}
\begin{document}
\section*{UnsolvedMath ID: 20003134}
\textbf{Source problem code:} catalogue entry
\subsection{Definitions and mathematical statement}
The source asks for an appropriate notion of Morse function on an Alexandrov space and examples; it does not prescribe a unique definition. The supplied aid proposes a flow-Morse class, distinct from asserting that ordinary smooth Morse theory already applies. In this proposed class a continuous locally semiconcave function on a compact finite-dimensional Alexandrov space has isolated metric-gradient critical points, regular bands admit compatible product structures, and sufficiently small passages through critical levels have stable local attachment pairs. A metric-gradient critical point means $\nabla f=0$ in the tangent-cone gradient sense. A regular-band product identifies $f^{-1}([a,b])$ with a level set times $[a,b]$ when no critical value is crossed; stable local pairs record the local sublevel attachment up to the chosen equivalence for all sufficiently small level changes. The exact full axioms and equivalence of pairs are not given in the record, so this is a proposed framework, not a claimed canonical definition or general Morse theorem.

The explicit family uses a nonempty compact Alexandrov space $\Sigma$ of curvature at least one in the spherical-suspension setting. With the usual diameter convention, for example $\operatorname{diam}\Sigma\leq\pi$, form
\[
X=\operatorname{Susp}(\Sigma)=([0,\pi]\times\Sigma)/(\{0\}\times\Sigma,\{\pi\}\times\Sigma\text{ each collapsed}),
\qquad f([t,\xi])=\cos t,
\]
with the spherical warped metric $\dd t^2+\sin^2t\,\dd s_\Sigma^2$. The aid asserts $f$ is flow-Morse with exactly the north and south poles critical. On $0<t<\pi$ the gradient is $-\sin t\,\partial_t$ and its norm is $\sin t$, so every closed band away from the poles is the explicit product with $\Sigma$.

For sublevel sets, the minimum has local attachment pair $(C\Sigma,\varnothing)$ and the maximum has $(C\Sigma,\Sigma)$, where $C\Sigma$ is the topological cone and $\Sigma$ its base. The maximum contributes
\[
H_k(C\Sigma,\Sigma;\Z)\cong\widetilde H_{k-1}(\Sigma;\Z)\qquad(k\geq1).
\]
The minimum contributes the usual degree-zero copy of $\Z$. These stable pairs are directly visible from the suspension coordinates.

Taking $\Sigma=\R\mathrm P^2$ with its standard curvature-one quotient metric gives
\[
H_2(C\R\mathrm P^2,\R\mathrm P^2;\Z)\cong\Z/2\Z.
\]
Thus a critical attachment need not have the ordinary smooth-Morse relative homology of a single integer-index handle $(D^\lambda,S^{\lambda-1})$, which is a free $\Z$ in one degree. The obstruction is to using that integer index alone as the complete attachment model, not to every conceivable Morse-type theory with richer local link data. No perturbation, genericity, or classification theorem for all Alexandrov spaces is supplied.
\subsection{Short English statement}
The supplied proposal uses metric gradients, product regular bands, and local attachment pairs. Suspension height gives an explicit two-critical-point family whose maximum remembers the reduced homology of the suspended space. Suspending the projective plane produces torsion in the critical attachment, so an ordinary single-cell Morse index cannot capture all the local data. The general definition and theorem package still need to be fixed.
\subsection{Sources}
\begin{itemize}
\item Local registry: permanent catalogue entry, exact source code catalogue entry; \texttt{aim-topology-notes.json}, \texttt{notes[221]}.
\item AIM workshop \emph{Manifolds with nonnegative sectional curvature}, Problem 9, \url{https://aimath.org/WWN/nnsectcurvature/nnsectcurvature.pdf}. The extracted footer is dated 26 October 2007 and says compiled by M. Kerin; compilation is not a proposer attribution. The concatenated ``12'' is pagination/footer material.
\item Supplied research attempt 1, registry updated 14 August 2026: partial result, low novelty confidence, machine-generated and not peer reviewed. It supplies the proposed flow-Morse framework and suspension calculation, without complete authoritative references for gradient theory or a formal axiomatic definition.
\item Original remarks and literature field are empty. Verification reports that the official PDF contains the mathematical sentence but gives no further definition or theorem citation.
\end{itemize}
\end{document}
